\clearpage
\appendix
\section{Supplementary information}

\subsection{Proof of the graph theorem}
\label{supp:proof}

Initialize the supplied output cotangents and set all other node cotangents to
zero.  Visit nodes in reverse topological order.  Suppose $v$ is the next node.
Every child of $v$ occurs later in forward order and has already been visited,
so $\bar z_v$ is the sum of all downstream path contributions.  For each
parent $u$, execute
\begin{equation}
 \bar z_u\mathrel{+}=P_{vu}(\bar z_v;c_v,\theta_v,\mathrm{static}),
\end{equation}
and, when the local parameter derivative is requested, execute
\begin{equation}
 \bar\theta_v\mathrel{+}=Q_v(\bar z_v;c_v,\theta_v,\mathrm{static}).
\end{equation}
By trace sufficiency, these are exactly the selected local transpose actions.
Addition handles several children, repeated edges, fan-out, residual branches,
and shared parameters.  Induction over reverse order therefore gives the
selected reverse-mode VJP at every requested node and parameter.  If every
$P,Q$ is generated by gather, scatter, copy, sign change, fixed address shifts,
and addition, their sequential execution introduces no operation outside that
set.  A local fixed non-dyadic scale or other declared operation remains in the
global evaluator and must be reported. \hfill$\square$

\subsection{Composition and finite recurrence}

The proof does not require the global selected Jacobian to be a partial
permutation.  A residual sum has local Jacobian $[I\ I]$, so its transpose
duplicates a cotangent.  Several selected outputs may scatter into the same
input; overlapping convolution and parameter sharing sum contributions.  The
global selected Jacobian can consequently contain multiple nonzeros per row or
column and integer coefficients greater than one.

For a recurrent model, apply the theorem to the stated finite unrolling.  A
carry/input gate and any parameter-sign code must be retained at every time
step whose gradient is requested.  Parameters on a recurrent edge are shared,
so their contributions accumulate.  The theorem does not establish convergence
of an infinite-horizon process or equivalence after a route-changing state
update.

\subsection{Tie selection and active-set stability}

For affine candidate scores, lexicographic ordering by value and negative
candidate index selects the first winner.  At an exact tie this defines a
selected limiting branch derivative rather than a unique classical derivative.
Suppose the original winner and runner-up differ by $\gamma>0$.  If every score
perturbation has magnitude at most $\epsilon$, the perturbed winner is at least
$s_1-\epsilon$ and every competitor is at most $s_2+\epsilon$.  The remaining
gap is at least $\gamma-2\epsilon>0$, so the route is unchanged.  At
$\gamma=0$, no positive radius follows from this bound.

The implementation must compute the winner and forward value through one
canonical primitive.  For example, a pairwise maximum can split a tie in an
autograd rule while an independently recorded $\geq$ mask routes all cotangent
through one branch.  Random tests away from ties cannot detect that semantic
mismatch.

\subsection{Conditional information bound}

Hold all reverse-available state fixed.  Call two realized local executions
distinguishable when some admissible cotangent yields different requested
transpose actions.  An exact no-recomputation reverse evaluator maps its stored
code and the cotangent to one action.  If two distinguishable executions shared
a code, the evaluator would receive identical available inputs for both and
return the same action, a contradiction.  Its code is therefore injective over
$M$ distinguishable classes.  Counting gives the fixed-width lower bound
$\lceil\log_2 M\rceil$.  The source-coding inequality gives expected length at
least $H(C)$ for a uniquely decodable variable-length code under a stated
distribution.

The conditioning matters.  Re-executing the forward, reconstructing values
through an invertible map, exploiting joint correlations between nodes, or
storing a different sufficient statistic changes the available information and
can change the bound.  It is not a universal lower bound on training memory.

\subsection{Counterexamples and exclusions}

\paragraph{Dense parameter gradient.}
Let $y=Wx$ and choose $x\ne x'$ while holding graph structure, shapes, and
weights fixed.  For the same nonzero output cotangent $q$, the weight VJPs are
$qx^T$ and $q{x'}^T$.  A route-free certificate that omits $x$ cannot return
both.  Constraining $W$ to powers of two can simplify the state action but does
not eliminate this activation outer product.

\paragraph{Smooth mixture.}
In $y=\sum_i\alpha_i(x)v_i$, both mixture weights and their derivatives vary
continuously with $x$.  A top expert index does not determine the pullback
unless the mixture is made genuinely hard and the derivative through its gate
is explicitly selected or stopped.

\paragraph{Ordinary normalization.}
Layer normalization has a Jacobian that depends on the realized mean and
variance.  A finite winner index does not determine it.  A variant that freezes
a quantized exponent may retain that exponent as auxiliary state, while its
centering reduction and any non-dyadic width scale remain explicit.

\paragraph{Route-changing replay.}
A certificate recorded at one internal state selects one piecewise program.
If clamping, parameter mutation, or iterative settling crosses an active-set
boundary, replaying the old certificate evaluates the old selected program.
The margin proposition gives a sufficient local no-crossing condition, not a
global convergence result.

\paragraph{Whole-model claim.}
Embedding addresses, loss logits and targets, dense/smooth modules, optimizer
moments, RNG state not represented by masks, and scheduling state do not vanish
because a hard-routed trunk is certified.  The artifact must name the exact
requested VJP and account for each boundary.

\subsection{MorphSSM boundary-accounting sensitivity}

Table~\ref{tab:accounting-sensitivity} exposes three conventions using the
same raw measurement.  The headline draft ratio divides the registered
checkpoint retained set by payload bytes after a supplied output cotangent.
Adding portable metadata changes little.  Adding both metadata and loss logits
reduces the ratio, and the smallest cell falls below 2$\times$.  The final
five-arm benchmark must choose a single derivative boundary for all arms and
show this sensitivity rather than hiding it.

\begin{table}[ht]
\centering
\caption{Sensitivity of checkpoint/certificate retained-byte ratio.  ``Full
portable + logits'' is serialized certificate bytes plus the separately
measured loss-boundary logits.}
\label{tab:accounting-sensitivity}
\small
\begin{tabular}{rrrrrrrr}
\toprule
$d$ & $L$ & $T$ & payload B & serialized B & payload ratio & serialized ratio & full portable + logits ratio \\
\midrule
64  & 2 & 96  & 36,084  & 37,431  & 2.45$\times$ & 2.36$\times$ & 1.70$\times$ \\
128 & 4 & 96  & 87,588  & 89,799  & 2.97$\times$ & 2.90$\times$ & 2.49$\times$ \\
256 & 4 & 192 & 329,928 & 332,173 & 3.07$\times$ & 3.05$\times$ & 2.80$\times$ \\
\bottomrule
\end{tabular}
\end{table}

\subsection{Claim-to-evidence boundary}

\begin{table}[ht]
\centering
\caption{Claims permitted by the current draft evidence.}
\small
\begin{tabular}{p{0.29\linewidth}p{0.18\linewidth}p{0.45\linewidth}}
\toprule
Claim & Status & Required wording or missing gate \\
\midrule
Graph composition theorem & established & finite executed graph, declared selected derivative, supplied boundary state \\
Additive max/min state and parameter certificate & established & additive parameterization and deterministic tie contract \\
Exact reference replay & established in tested scope & 38 tests plus matched FP32 trajectories; not universal operator coverage \\
MorphSSM retained-state representation & measured, configuration-specific & 2.45--3.07$\times$ payload versus registered checkpoint accounting; show boundary sensitivity \\
Process/allocator peak advantage & open & direct peak measurement on matched five-arm cells \\
Memory-traffic advantage & open & measured or validated modeled bytes moved, including packing and decoding \\
Compiled full-cycle latency advantage & open & five-arm benchmark, consequential scale, exactness first \\
Generic activation-free training & rejected & continuous/loss/optimizer boundaries remain explicit \\
New gradient or learning rule & rejected & certificate replay is selected reverse-mode accumulation \\
Whole-model coverage & rejected & operator-specific VJP contract only \\
Energy or hardware advantage & outside this paper & separate hardware work with matched physical baselines \\
\bottomrule
\end{tabular}
\end{table}

\clearpage
\subsection{Submission artifact checklist}

\begin{enumerate}[leftmargin=1.5em]
  \item Move the reference primitives, packed schema, tests, and raw
  retained-state measurement into versioned TikiTorch paths.
  \item Bind every result to the submission commit, clean/dirty digest,
  environment lock, full command, device, seed, and input/model digests.
  \item Freeze the five-arm configurations, boundary convention, tolerances,
  scale ladder, repetitions, and decision thresholds before execution.
  \item Include eager and compiled autograd, checkpoint/recompute, unpacked
  compiled custom VJP, and packed compiled replay on the same graph.
  \item Measure full-cycle latency, peak memory, retained bytes, recomputation,
  and traffic; retain compiler fallbacks and failed/OOM rows.
  \item Regenerate all tables from append-only raw records with one
  deterministic reducer and verify the clean-room reproduction command.
\end{enumerate}
